% Options for packages loaded elsewhere
% Options for packages loaded elsewhere
\PassOptionsToPackage{unicode}{hyperref}
\PassOptionsToPackage{hyphens}{url}
\PassOptionsToPackage{dvipsnames,svgnames,x11names}{xcolor}
%
\documentclass[
  german,
  10pt,
  twoside]{article}
\usepackage{xcolor}
\usepackage[top=2.5cm,bottom=2.5cm,left=2cm,right=2cm]{geometry}
\usepackage{amsmath,amssymb}
\setcounter{secnumdepth}{5}
\usepackage{iftex}
\ifPDFTeX
  \usepackage[T1]{fontenc}
  \usepackage[utf8]{inputenc}
  \usepackage{textcomp} % provide euro and other symbols
\else % if luatex or xetex
  \usepackage{unicode-math} % this also loads fontspec
  \defaultfontfeatures{Scale=MatchLowercase}
  \defaultfontfeatures[\rmfamily]{Ligatures=TeX,Scale=1}
\fi
\usepackage{lmodern}
\ifPDFTeX\else
  % xetex/luatex font selection
\fi
% Use upquote if available, for straight quotes in verbatim environments
\IfFileExists{upquote.sty}{\usepackage{upquote}}{}
\IfFileExists{microtype.sty}{% use microtype if available
  \usepackage[]{microtype}
  \UseMicrotypeSet[protrusion]{basicmath} % disable protrusion for tt fonts
}{}
\usepackage{setspace}
\makeatletter
\@ifundefined{KOMAClassName}{% if non-KOMA class
  \IfFileExists{parskip.sty}{%
    \usepackage{parskip}
  }{% else
    \setlength{\parindent}{0pt}
    \setlength{\parskip}{6pt plus 2pt minus 1pt}}
}{% if KOMA class
  \KOMAoptions{parskip=half}}
\makeatother
% Make \paragraph and \subparagraph free-standing
\makeatletter
\ifx\paragraph\undefined\else
  \let\oldparagraph\paragraph
  \renewcommand{\paragraph}{
    \@ifstar
      \xxxParagraphStar
      \xxxParagraphNoStar
  }
  \newcommand{\xxxParagraphStar}[1]{\oldparagraph*{#1}\mbox{}}
  \newcommand{\xxxParagraphNoStar}[1]{\oldparagraph{#1}\mbox{}}
\fi
\ifx\subparagraph\undefined\else
  \let\oldsubparagraph\subparagraph
  \renewcommand{\subparagraph}{
    \@ifstar
      \xxxSubParagraphStar
      \xxxSubParagraphNoStar
  }
  \newcommand{\xxxSubParagraphStar}[1]{\oldsubparagraph*{#1}\mbox{}}
  \newcommand{\xxxSubParagraphNoStar}[1]{\oldsubparagraph{#1}\mbox{}}
\fi
\makeatother


\usepackage{graphicx}
\makeatletter
\newsavebox\pandoc@box
\newcommand*\pandocbounded[1]{% scales image to fit in text height/width
  \sbox\pandoc@box{#1}%
  \Gscale@div\@tempa{\textheight}{\dimexpr\ht\pandoc@box+\dp\pandoc@box\relax}%
  \Gscale@div\@tempb{\linewidth}{\wd\pandoc@box}%
  \ifdim\@tempb\p@<\@tempa\p@\let\@tempa\@tempb\fi% select the smaller of both
  \ifdim\@tempa\p@<\p@\scalebox{\@tempa}{\usebox\pandoc@box}%
  \else\usebox{\pandoc@box}%
  \fi%
}
% Set default figure placement to htbp
\def\fps@figure{htbp}
\makeatother



\ifLuaTeX
\usepackage[bidi=basic]{babel}
\else
\usepackage[bidi=default]{babel}
\fi
% get rid of language-specific shorthands (see #6817):
\let\LanguageShortHands\languageshorthands
\def\languageshorthands#1{}
\ifLuaTeX
  \usepackage[german]{selnolig} % disable illegal ligatures
\fi


\setlength{\emergencystretch}{3em} % prevent overfull lines

\providecommand{\tightlist}{%
  \setlength{\itemsep}{0pt}\setlength{\parskip}{0pt}}



 
\usepackage[style=authoryear,backend=biber,style=authoryear-comp,uniquename=false,uniquelist=false]{biblatex}
\addbibresource{../references.bib}


\usepackage{booktabs}
\usepackage{longtable}
\usepackage{array}
\usepackage{multirow}
\usepackage{wrapfig}
\usepackage{float}
\usepackage{colortbl}
\usepackage{pdflscape}
\usepackage{tabu}
\usepackage{threeparttable}
\usepackage{threeparttablex}
\usepackage[normalem]{ulem}
\usepackage{makecell}
\usepackage{xcolor}
\usepackage{csquotes}
\usepackage{booktabs}
\usepackage{array}
\usepackage{multirow}
\usepackage{wrapfig}
\usepackage{float}
\usepackage{colortbl}
\usepackage{pdflscape}
\usepackage{threeparttable}
\usepackage{threeparttablex}
\usepackage[normalem]{ulem}
\usepackage{makecell}
\usepackage{xcolor}
\usepackage{multicol}
\usepackage{fancyhdr}
\usepackage{titlesec}
\usepackage{etoolbox}
\usepackage{tabularx}
\usepackage{afterpage}
\usepackage{tikz}
\usepackage{longtable}



% Suppress default maketitle
\renewcommand{\maketitle}{}

% Define WISTA colors (Havelock Blue as main)
\definecolor{wistaBlue}{RGB}{65,125,180}
\definecolor{wistaLightBlue}{RGB}{120,170,215}
\definecolor{wistaDarkGray}{RGB}{64,64,64}
\definecolor{wistaLightGray}{RGB}{128,128,128}

% Section formatting (##) - number 20% bigger, lines closer, text closer to bottom line
% e.g., "2" with lines around "Einleitung"
\titleformat{\section}[display]
  {\normalfont\large\bfseries}
  {\color{wistaBlue}\Large\bfseries\thesection\\\noindent\rule{\columnwidth}{1.5pt}}
  {0pt}
  {\vspace{0pt}\color{wistaDarkGray}\large\bfseries}
  [\vspace{0pt}\noindent{\color{wistaBlue}\rule{\columnwidth}{1.5pt}}\vspace{0.05cm}]
\titlespacing*{\section}{0pt}{1.5ex plus 0.3ex minus .1ex}{0.2ex plus .1ex}

% Subsection formatting (###) - grey text with blue line under
% e.g., "2.1 Handlungsbedarf"
\renewcommand{\thesubsection}{\thesection.\arabic{subsection}}
\titleformat{\subsection}
  {\color{wistaDarkGray}\normalsize\bfseries}
  {\color{wistaDarkGray}\normalsize\bfseries\thesubsection}
  {0.5em}
  {\color{wistaDarkGray}\normalsize\bfseries}
  [\vspace{0pt}\noindent{\color{wistaBlue}\rule{\columnwidth}{0.8pt}}\vspace{0.05cm}]
\titlespacing*{\subsection}{0pt}{1ex plus 0.2ex minus .1ex}{0.2ex plus .1ex}
  
% Subsubsection formatting (####) - grey text, no underline
% e.g., "2.1.1 Steigende Datenmenge"
\renewcommand{\thesubsubsection}{\thesubsection.\arabic{subsubsection}}
\titleformat{\subsubsection}
  {\color{wistaDarkGray}\normalsize\bfseries}
  {\color{wistaDarkGray}\normalsize\bfseries\thesubsubsection}
  {0.5em}
  {\color{wistaDarkGray}\normalsize\bfseries}

% Paragraph formatting (4th level) - numbered as subsection.subsubsection.paragraph
\renewcommand{\theparagraph}{\thesubsubsection.\arabic{paragraph}}
\titleformat{\paragraph}[runin]
  {\color{wistaDarkGray}\normalsize\bfseries}
  {\color{wistaBlue}\normalsize\bfseries\theparagraph}
  {0.5em}
  {\color{wistaDarkGray}\normalsize\bfseries}
  [.]
\titlespacing*{\paragraph}{0pt}{1ex plus 0.5ex minus 0.2ex}{0.5em}

% Header and footer setup with blue line under header
\pagestyle{fancy}
\fancyhf{}
\fancyhead[L]{\small\textit{Oliver Hauke, Annette Kristiansen}}
\fancyhead[R]{\small\textit{Effiziente Analyse großer Forschungsdaten}}
\fancyfoot[L]{\small\thepage}
\fancyfoot[R]{\small Statistisches Bundesamt | WISTA | {\color{wistaBlue}6} | 2025}
\renewcommand{\headrulewidth}{0pt}
\renewcommand{\footrulewidth}{0pt}
% Add blue line under header
\renewcommand{\headrule}{\color{wistaBlue}\hrule width\headwidth height 1pt \vskip 2pt}
\renewcommand{\headrulewidth}{1pt}

% Make all figures span both columns (full width)
\let\origfigure\figure
\let\endorigfigure\endfigure
\renewenvironment{figure}[1][htbp]{%
  \origfigure*[#1]%
}{%
  \endorigfigure*%
}

% Allow manual control of table* environment for full-width tables
% Tables by default stay in single column, but can use table* manually



% Make bibliography full-width (one column) - for biblatex
\defbibheading{bibliography}[\bibname]{%
  \onecolumn
  \section*{#1}%
  \markboth{#1}{#1}%
}



% Footnotes stay within column
\usepackage[stable]{footmisc}

% Custom footnote formatting: "|^1" in dark WISTA gray
\renewcommand{\thefootnote}{{\color{wistaDarkGray}|$^{\arabic{footnote}}$}}

% Pifont for dingbat symbols
\usepackage{pifont}

% Better table handling for two-column mode
\usepackage{array,tabularx,calc}

% Custom WISTA table format
\usepackage{caption}
\usepackage{xcolor}
\usepackage{colortbl}
\usepackage{array}

% Define custom caption format for WISTA tables
\DeclareCaptionFormat{wista}{%
  {\color{wistaLightBlue}\bfseries #1#2}\par\vspace{0.2em}%
  {\color{wistaDarkGray}\normalfont #3}\par\vspace{0.3em}%
}

% Set up table captions with custom WISTA format
\captionsetup[table]{
  position=top,
  format=wista,
  justification=justified,
  singlelinecheck=false,
  labelsep=none,
  skip=0.3em
}

% Better table spacing and white inner lines
\setlength{\tabcolsep}{6pt}
\renewcommand{\arraystretch}{1.15}

% Set white lines for tables by default, blue for first column
\AtBeginDocument{
  \arrayrulecolor{white}
  \setlength{\arrayrulewidth}{0.3pt}
}

% Define commands for colored table borders
\newcommand{\bluevrule}{\color{wistaBlue}\vrule}
\newcommand{\whitevrule}{\color{white}\vrule}

% Define WISTA table colors
\definecolor{wistaTableBg}{RGB}{240,248,255}     % Very light blue background
\definecolor{wistaHeaderBg}{RGB}{220,230,240}    % Light blue header background
\definecolor{wistaFirstCol}{RGB}{65,125,180}     % Blue for first column lines

% Custom WISTA table command
\newcommand{\wistaTableStart}{%
  \rowcolors{2}{wistaTableBg}{wistaTableBg}%
  \arrayrulecolor{wistaFirstCol}%
}

% Custom commands for table styling
\newcommand{\wistaHeaderRow}[1]{%
  \rowcolor{wistaHeaderBg}%
  \arrayrulecolor{wistaDarkGray}%
  #1 \\
  \arrayrulecolor{white}%
}

\newcommand{\wistaFirstColRule}{\arrayrulecolor{wistaFirstCol}}
\newcommand{\wistaWhiteRule}{\arrayrulecolor{white}}

% Custom bullet points with > character - 35% width, 95% height, black extra bold, half line spacing
\usepackage{enumitem}
\setlist[itemize]{label={\scalebox{0.35}[0.95]{\fontseries{bx}\selectfont >}}, leftmargin=1.5em, itemsep=5pt, topsep=3pt, parsep=0pt}

% Redefine ref to include pifont 247 (arrow) for figure/table/section references
\let\oldref\ref
\renewcommand{\ref}[1]{\oldref{#1}\raisebox{-0.2ex}{\scriptsize\color{wistaBlue}\ding{247}}}
\makeatletter
\@ifpackageloaded{caption}{}{\usepackage{caption}}
\AtBeginDocument{%
\ifdefined\contentsname
  \renewcommand*\contentsname{Inhaltsverzeichnis}
\else
  \newcommand\contentsname{Inhaltsverzeichnis}
\fi
\ifdefined\listfigurename
  \renewcommand*\listfigurename{Abbildungsverzeichnis}
\else
  \newcommand\listfigurename{Abbildungsverzeichnis}
\fi
\ifdefined\listtablename
  \renewcommand*\listtablename{Tabellenverzeichnis}
\else
  \newcommand\listtablename{Tabellenverzeichnis}
\fi
\ifdefined\figurename
  \renewcommand*\figurename{Abbildung}
\else
  \newcommand\figurename{Abbildung}
\fi
\ifdefined\tablename
  \renewcommand*\tablename{Tabelle}
\else
  \newcommand\tablename{Tabelle}
\fi
}
\@ifpackageloaded{float}{}{\usepackage{float}}
\floatstyle{ruled}
\@ifundefined{c@chapter}{\newfloat{codelisting}{h}{lop}}{\newfloat{codelisting}{h}{lop}[chapter]}
\floatname{codelisting}{Listing}
\newcommand*\listoflistings{\listof{codelisting}{Listingverzeichnis}}
\makeatother
\makeatletter
\makeatother
\makeatletter
\@ifpackageloaded{caption}{}{\usepackage{caption}}
\@ifpackageloaded{subcaption}{}{\usepackage{subcaption}}
\makeatother
\usepackage{bookmark}
\IfFileExists{xurl.sty}{\usepackage{xurl}}{} % add URL line breaks if available
\urlstyle{same}
\hypersetup{
  pdftitle={EFFIZIENTE AUSWERTUNG DES BUSINESS-TAX-PANELS MITHILFE VON R},
  pdfauthor={Oliver Hauke; Annette Kristiansen},
  pdflang={de},
  colorlinks=true,
  linkcolor={blue},
  filecolor={Maroon},
  citecolor={Blue},
  urlcolor={Blue},
  pdfcreator={LaTeX via pandoc}}


\title{EFFIZIENTE AUSWERTUNG DES BUSINESS-TAX-PANELS MITHILFE VON R}
\author{Oliver Hauke \and Annette Kristiansen}
\date{2025-12-04}
\begin{document}
\maketitle

% Custom WISTA-style title page
\begin{titlepage}
\thispagestyle{empty}
\pagestyle{empty}
\onecolumn

% Horizontal blue line at top (header line)
\noindent{\color{wistaBlue}\rule{\textwidth}{1pt}}

\vspace{-0.5cm}

% Vertical blue line on the left (6.3cm from left edge, starts below header line, stops above footer)
\begin{tikzpicture}[remember picture, overlay]
  \draw[wistaBlue, line width=2pt] 
    ([xshift=6.3cm, yshift=-4cm]current page.north west) -- 
    ([xshift=6.3cm, yshift=3.5cm]current page.south west);
\end{tikzpicture}

\vspace*{0.5cm}

% Left side - short author bios (right-aligned to the vertical line)
\begin{minipage}[t]{3.9cm}
\selectlanguage{ngerman}%
\raggedleft
\hyphenpenalty=0
\tolerance=1000
\fontsize{7}{8.5}\selectfont
\vspace{0pt}%
{\fontsize{8}{9.5}\selectfont\color{wistaBlue}\textbf{Oliver Hauke}}\\[0.15cm]
ist Mathematiker und IT-Referent im Kompetenzzentrum „Analyse und Auswertung" des Statistischen Bundesamtes. Er verantwortet die Weiterentwicklung der technischen Infrastruktur des Forschungsdatenzentrums und berät das Netzwerk empirische Steuerforschung in der effizienten Nutzung der Systeme.\\[0.5cm]

{\fontsize{8}{9.5}\selectfont\color{wistaBlue}\textbf{Annette Kristiansen}}\\[0.15cm]
ist Ökonomin und Referentin im Referat „Unternehmenssteuern" des Statistischen Bundesamtes. Sie beschäftigt sich mit der Zusammenführung der Unternehmenssteuerstatistiken und betreut das Netzwerk empirische Steuerforschung. Für ihre externe Promotion an der Universität Bremen untersucht sie die technologische Vielfalt multinationaler Unternehmen.

\end{minipage}%
\hspace{0.4cm}% Gap to vertical line (left side)
\hspace{0.4cm}% Gap from vertical line (right side)
% Right side - title, authors, keywords, abstracts
\begin{minipage}[t]{11.8cm}
\raggedright
\vspace{1.2cm}%

% Title (uppercase via command, not bold)
{\LARGE\color{wistaDarkGray}\begin{spacing}{1.3}\MakeUppercase{Effiziente Auswertung des Business-Tax-Panels mithilfe von R}\end{spacing}}

\vspace{-0.2cm}

% Blue horizontal line (underline for title)
{\color{wistaBlue}\rule{\linewidth}{0.5pt}}

\vspace{0.5cm}

% Authors inline, dark gray
{\large\color{wistaDarkGray} Oliver Hauke, Annette Kristiansen}

\vspace{0.05cm}

% Blue line (underline for authors)
{\color{wistaBlue}\rule{\linewidth}{0.5pt}}

\vspace{0.5cm}

% Keywords (pifont 247 arrow then bold label and blue text)
\hangindent=1.2em\hangafter=1
{\color{wistaBlue}\ding{247}\ {\bfseries Schlüsselwörter:} Effiziente Analysen, R-Programmierung, Forschungsdaten, Unternehmenssteuerstatistik, Business-Tax-Panel, Netzwerk empirische Steuerforschung}

\vspace{0.3cm}

% Zusammenfassung (uppercase via command)
{\color{wistaDarkGray}\bfseries\MakeUppercase{Zusammenfassung}}\\[0.15cm]
\small
Das Forschungsdatenzentrum des Statistischen Bundesamtes baut seine technische Infrastruktur gezielt aus, um der Wissenschaft erweiterte Analysemöglichkeiten amtlicher Daten bereitzustellen. Seit November 2025 steht Forschenden exklusiver Zugang zu diesen erweiterten Systemressourcen in R und Stata zur Verfügung. R-basierte Tools -- insbesondere Apache Arrow, Parquet und DuckDB -- ermöglichen es, gängige Analysen großer Forschungsdatensätze in Echtzeit auszuführen. Am Beispiel des Business-Tax-Panel (2013 bis 2019) werden die erzielbaren Effizienzgewinne bei der Auswertung umfangreicher steuerstatistischer Daten im Rahmen des Netzwerks empirische Steuerforschung demonstriert.

\vspace{0.5cm}

% English keywords and abstract (italicized)
\hangindent=1.2em\hangafter=1
{\itshape\color{wistaBlue}\ding{247}\ {\bfseries Keywords:} Efficient analysis -- R programming -- research data -- corporate tax statistics -- Business Tax Panel -- empirical tax research network}

\vspace{0.3cm}

{\itshape\color{wistaDarkGray}\bfseries\MakeUppercase{Abstract}}\\[0.15cm]
\small\itshape
The Research Data Centre at the Federal Statistical Office is strategically expanding its technical infrastructure to provide researchers with enhanced analytical capabilities for official data. Since November 2025, researchers have had exclusive access to these extended system resources in R and Stata. R-based tools -- particularly Apache Arrow, Parquet, and DuckDB -- enable real-time analysis of large research datasets. Using the Business Tax Panel (2013--2019), this paper demonstrates the achievable efficiency gains in analysing large-scale tax statistics within the empirical tax research network.

\end{minipage}

\vfill

\vspace{0.5cm}
% Footer matching document style
\noindent\makebox[\textwidth]{%
  \small 1%
  \hfill%
  \small Statistisches Bundesamt | WISTA | {\color{wistaBlue}6} | 2025%
}

\end{titlepage}

% Stay in one-column mode for TOC (will be switched later)
\pagestyle{fancy}

\renewcommand*\contentsname{Inhaltsverzeichnis}
{
\hypersetup{linkcolor=}
\setcounter{tocdepth}{3}
\tableofcontents
}

\setstretch{1}
\section{Anwendungsfall Business-Tax-Panel}\label{sec-btp}

\subsection{Datengrundlage}\label{datengrundlage}

\begin{itemize}
\tightlist
\item
  Umfang über verfügbaren Arbeitsspeicher
\item
  Besonders breit mit 3000 Merkmalen, 80 Mio. Obs.
\end{itemize}

\subsection{Kennzahlen}\label{kennzahlen}

\subsection{Auswertungsbedarfe}\label{auswertungsbedarfe}

Breites Spektrum: Datenaufbereitung/-bereitstellung und Auswertung von
Mikrodaten durch die Wissenschaft.

\begin{enumerate}
\def\labelenumi{\arabic{enumi}.}
\item
  deskriptive Eckdaten, für BTP als Panel im Quer- und Längsschnitt:
  Fallzahl pro Statistik und Jahr (s.MDR Produkt S. 9f
  https://www.forschungsdatenzentrum.de/sites/default/files/btp\_2013-2019\_on-site\_mdr-prod.pdf)
\item
  Verlustvorträge (s. Vortrag NeSt) als Indikator für \ldots{} und
  Einflussgrößen auf Verlustvorträge, für Vermutung von Einfluss der
  Unternehmensgröße (großes Gesamtvolumen, Löhne, Spenden, tätige
  Personen; Auslandskontrolle oder WZ) auf positive Verlustvorträge.
\end{enumerate}

Aus technischer Sicht, ist die Kernaufgabe den Großteil (von 2991
Variablen) effizient zu entfernen, sodass lediglich der relevante
Ausschnitt von kleinem Umfang ausgewertet werden kann.

Beide Anwendungsfälle können bereits im mit den folgenden X von 3000
Variablen untersucht werden:

\begin{table*}

\caption{\label{tbl-variables}Betrachtete Variablen des BTP}

\centering{

\centering\begingroup\fontsize{9}{11}\selectfont


\begin{tabular}{|>{}l|||>{}l|||>{}l|}
\hline
\cellcolor[HTML]{dce6f0}{\textcolor[HTML]{404040}{\textbf{Variable}}} & \cellcolor[HTML]{dce6f0}{\textcolor[HTML]{404040}{\textbf{Statistik}}} & \cellcolor[HTML]{dce6f0}{\textcolor[HTML]{404040}{\textbf{Beschreibung}}}\\
\hline
\cellcolor{white}{\textcolor{black}{jahr}} & \cellcolor[HTML]{f0f8ff}{Allgemein} & \cellcolor[HTML]{f0f8ff}{Berichtsjahr}\\
\cellcolor{white}{\textcolor{black}{verk}} & \cellcolor[HTML]{f0f8ff}{Allgemein} & \cellcolor[HTML]{f0f8ff}{}\\
\cellcolor{white}{\textcolor{black}{k\_k65270}} & \cellcolor[HTML]{f0f8ff}{KStSt} & \cellcolor[HTML]{f0f8ff}{Verlustvorträge}\\
\cellcolor{white}{\textcolor{black}{k\_k65823}} & \cellcolor[HTML]{f0f8ff}{KStSt} & \cellcolor[HTML]{f0f8ff}{Gesamtbetrag der Einkünfte}\\
\cellcolor{white}{\textcolor{black}{k\_c15018}} & \cellcolor[HTML]{f0f8ff}{KStSt} & \cellcolor[HTML]{f0f8ff}{Summe der Umsätze/Löhne/Gehälter}\\
\cellcolor{white}{\textcolor{black}{k\_ef20}} & \cellcolor[HTML]{f0f8ff}{KStSt} & \cellcolor[HTML]{f0f8ff}{WZ Generierung}\\
\cellcolor{white}{\textcolor{black}{k\_k65172}} & \cellcolor[HTML]{f0f8ff}{KStSt} & \cellcolor[HTML]{f0f8ff}{Spende}\\
\cellcolor{white}{\textcolor{black}{urs\_we\_tp\_stichtag}} & \cellcolor[HTML]{f0f8ff}{URS} & \cellcolor[HTML]{f0f8ff}{Tätige Personen}\\
\cellcolor{white}{\textcolor{black}{urs\_rt\_gruppen\_kennz}} & \cellcolor[HTML]{f0f8ff}{URS} & \cellcolor[HTML]{f0f8ff}{Statuskennzeichen, = 3/6 heißt "auslandskontrolliert"}\\
\hline
\end{tabular}
\endgroup{}

}

\end{table*}%

Kennzahlen um inhaltlose Testdaten zu generieren werden folgende
Eckdaten benötigt und im nächsten BTP berücksichtigt:

\begin{itemize}
\tightlist
\item
  Befüllung
\item
  Statistische Eckdaten univariate (Quantile), multivariate (gemeinsame
  Verteilung)
\end{itemize}

Update BTP, flexibel reagieren, neue MDR, DSB, \ldots{}

\begin{itemize}
\tightlist
\item
  Datenbereitstellung: MDR
\item
  Forschungsprojekt

  \begin{itemize}
  \tightlist
  \item
    Lorenzkurve -\textgreater{} Verteilung der Verlustvorträge
  \item
    Logistische Regression -\textgreater{} Einflussfaktoren auf
    Verlustvorträge
  \end{itemize}
\end{itemize}

\subsection{\ldots{}}\label{section}

\begin{itemize}
\tightlist
\item
  TODO: Zufügen von öffentlich verfügbaren Informationen zu Struktur und
  Verteilung des BTP. Verteilung über Statistik und Jahr, \ldots{}

  \begin{itemize}
  \tightlist
  \item
    was ist wichtig? Randverteilung, missings, vollständigkeit der
    variablen
  \item
    erstmal eine Tabelle mit Variablenanzahl, Entfernung durch Hinweis
    feld. Tabelle: Variablenzahl komplett, Variablenanzahl ohne
    Hinweisfeld, Variablenanzahl Befüllt?
  \item
    Hinweis Befüllungsgrad tlw. gering\ldots{} x Variablen fehlen zu
    90\%, 80\%, 70\%, 60\%, 50\% ?
  \item
    Verteilung der Zeilen auf Statistik und Jahr, inkl. Trend
  \item
    Verteilung auf Kombinationen genauer
  \end{itemize}
\item
  Anwendungsfälle volles Spektrum

  \begin{itemize}
  \tightlist
  \item
    Datenbereitstellung, interne Analysen, Schätzungen
  \item
    Auswertung durch Wissenschaft im FDZ
  \end{itemize}
\item
  Technische Herausforderung:

  \begin{itemize}
  \tightlist
  \item
    effiziente und verständliche Speicherung
  \item
    Herausfiltern großer Datenmengen irrelevanter Merkmale
  \end{itemize}
\end{itemize}

\textless!--

\subsection{Annette's Input}\label{annettes-input}

\emph{Das Business-Tax-Panel kombiniert die Angaben aus verschiedenen
Ertrags- sowie Umsatzsteuerstatistiken im Quer- und Längsschnitt.
Verfügbar sind Informationen aus den folgenden Statistiken :
Gewerbesteuer, Körperschaftsteuer, Umsatzsteuer-Voranmeldung und
-Veranlagung, Statistik über die Personengesellschaften und
Gemeinschaften, Einnahmenüberschussrechnung und einige Merkmale aus dem
Statistischen Unternehmensregister. Alle verwendeten Steuerarte sind ab
dem Berichtsjahr 2013 jährlich verfügbar, dies bildet den Beginn des
Beobachtungszeitraums. Aufgrund der Festsetzungsfrist von vier Jahren
bei den Veranlagungsstatistiken endet der Beobachtungszeitraum aktuell
mit dem Berichtsjahr 2020.}

\emph{Bei den verwendeten Statistiken handelt es sich um Vollerhebungen,
teilweise mit Erfassungsgrenze. Insgesamt liegen 77,8 Mio. Beobachtungen
vor, die auf 16,7 Mio. Einheiten zurückzuführen sind. Der Merkmalsumfang
variiert je nach Statistik zwischen 120 und 1.300 Merkmalen was
insgesamt zu einem Merkmalsumfang von über 2.700 Merkmalen führt.
Hierdurch sind Analyse von Anpassungsreaktionen sowie
Verteilungsanalysen für verschiedene Beobachtungszeiträume möglich.
Insgesamt deckt der Forschungsdatensatz eine Vielzahl an
Analysemöglichkeiten, die für die evidenzbasierte steuerpolitische
Entscheidungen nötig sind, ab.}

\emph{Für die vielfältigen Analysewünsche der Forschenden, den
Auswertungsbedarf des Fachbereiches bspw. für Sonderauswertungen und der
Erstellung des Metadatenreports des Forschungsdatenzentrums ist der
Forschungsdatensatz konzipiert. Für die Abdeckung aller Anforderungen
ist der Datensatz möglichst wenig beschränkt worden und dadurch
ressourcenintensive. Der damit einhergehende hohe Bedarf an
Speicherkapazität, Arbeitsspeicher und Prozessorleistung stellt sowohl
im Fachbereich als auch im Forschungsdatenzentrum eine Herausforderung
dar.}

==Wodurch \ldots==

\subsection{Überblick zum
Business-Tax-Panel}\label{uxfcberblick-zum-business-tax-panel}

Das Business-Tax-Panel (BTP) ist ein zentraler Forschungsdatensatz des
Statistischen Bundesamtes, der umfassende Informationen zur
Unternehmensbesteuerung in Deutschland bereitstellt. Das Panel
kombiniert Daten aus verschiedenen Steuerstatistiken miteinander und
ermöglicht somit eine umfassende Analyse der Besteuerung von Unternehmen
über Zeit.

Im BTP werden folgende Datensätze zusammengeführt:

\begin{itemize}
\tightlist
\item
  \textbf{Gewerbesteuerstatistik} (EVAS-Nr. 73511)
\item
  \textbf{Körperschaftsteuerstatistik} (EVAS-Nr. 73211)
\item
  \textbf{Umsatzsteuerstatistik - Voranmeldungen} (EVAS-Nr. 73311)
\item
  \textbf{Statistik über Personengesellschaften und Gemeinschaften}
  (EVAS-Nr. 73121)
\item
  \textbf{Umsatzsteuerstatistik - Veranlagungen} (EVAS-Nr. 73321)
\item
  \textbf{Einnahmenüberschussrechnung} aus der Lohn- und
  Einkommensteuerstatistik (EVAS-Nr. 73111)
\item
  Merkmale aus dem \textbf{Statistischen Unternehmensregister} (EVAS-Nr.
  52111)
\end{itemize}

\subsection{Datenstruktur und Umfang}\label{datenstruktur-und-umfang}

\textbf{Dimensionen des Datensatzes.}

Das BTP umfasst einen außergewöhnlich umfangreichen Datensatz mit
beeindruckenden Dimensionen:

\begin{table*}

\caption{\label{tbl-btp-stats}Übersicht über die Dimensionen des Business-Tax-Panel}

\centering{

\centering\begingroup\fontsize{10}{12}\selectfont


\begin{tabular}{|>{}l|||>{}l|}
\hline
\cellcolor[HTML]{dce6f0}{\textcolor[HTML]{404040}{\textbf{Merkmal}}} & \cellcolor[HTML]{dce6f0}{\textcolor[HTML]{404040}{\textbf{Wert}}}\\
\hline
\cellcolor{white}{\textcolor{black}{Beobachtungszeitraum}} & \cellcolor[HTML]{f0f8ff}{2013-2019 (jährliche Erweiterung)}\\
\cellcolor{white}{\textcolor{black}{Anzahl Einheiten (2013-2019)}} & \cellcolor[HTML]{f0f8ff}{>66 Millionen}\\
\cellcolor{white}{\textcolor{black}{Anzahl Merkmale}} & \cellcolor[HTML]{f0f8ff}{>2.700}\\
\cellcolor{white}{\textcolor{black}{Erhebungsart}} & \cellcolor[HTML]{f0f8ff}{Vollerhebung}\\
\cellcolor{white}{\textcolor{black}{Datenquelle}} & \cellcolor[HTML]{f0f8ff}{Sekundärerhebung (Finanzverwaltung)}\\
\hline
\end{tabular}
\endgroup{}

}

\end{table*}%

\textbf{Hauptmerkmale des Datensatzes.}

Aus den verschiedenen Steuervordrucken und darin verwendeten Kennzahlen
werden über 2.700 Merkmale erfasst, die verschiedene Aspekte der
Unternehmensbesteuerung abbilden:

\begin{itemize}
\tightlist
\item
  \textbf{Steuerliche Kennzahlen:} Bemessungsgrundlagen, Steuerbeträge,
  Vorsteuern aus verschiedenen Steuerarten
\item
  \textbf{Wirtschaftliche Kennzahlen:} Umsätze, Gewinne,
  Betriebsausgaben
\item
  \textbf{Strukturmerkmale:} Rechtsform, Wirtschaftszweig, Größenklassen
\item
  \textbf{Geografische Informationen:} Bundesland, Gemeindekennziffer
\item
  \textbf{Panelstruktur:} Längsschnittinformationen ermöglichen
  Zeitreihenanalysen
\end{itemize}

\textbf{Datenerhebung und Rechtsgrundlage.}

Die im BTP enthaltenen Steuerstatistiken sind Bundesstatistiken, die als
\textbf{Vollerhebung} durchgeführt werden. Es handelt sich um
\textbf{Sekundärerhebungen}: Die Erhebungsmerkmale werden aus den
Voranmeldungs- und Veranlagungsbescheiden von der Finanzverwaltung
entnommen und über deren Rechenzentren an die Statistischen Ämter der
Länder übermittelt.

Rechtsgrundlage ist das
\href{https://www.gesetze-im-internet.de/ststatg_1995/}{Gesetz über
Steuerstatistiken (StStatG)} vom 11. Oktober 1995.

\subsection{Technische Requirements}\label{technische-requirements}

\begin{itemize}
\tightlist
\item
  Größe in Datentypen
\item
  Vorbereitung OSS
\end{itemize}

\subsection{Forschungspotenzial und
Nutzung}\label{forschungspotenzial-und-nutzung}

\textbf{Typische Forschungsfragen.}

Das BTP wird für eine Vielzahl von Forschungsfragen genutzt:

\begin{enumerate}
\def\labelenumi{\arabic{enumi}.}
\tightlist
\item
  \textbf{Steuerliche Analysen:} Effekte von Steuerreformen auf
  Unternehmensentscheidungen
\item
  \textbf{Unternehmensdynamik:} Gründungen, Wachstum und Marktaustritte
\item
  \textbf{Wirtschaftsstruktur:} Branchenentwicklung und regionale
  Unterschiede
\item
  \textbf{Krisenanalysen:} Auswirkungen von wirtschaftlichen Schocks auf
  Unternehmen
\end{enumerate}

\textbf{Herausforderungen bei der Datenauswertung.}

Die Größe und Komplexität des BTP stellt besondere Anforderungen an die
Auswertungsinfrastruktur:

\begin{itemize}
\tightlist
\item
  \textbf{Rechenleistung:} Analysen über den gesamten Datensatz
  erfordern erhebliche Rechenkapazitäten
\item
  \textbf{Speicherbedarf:} Das Vorhalten des Datensatzes im
  Arbeitsspeicher ist oft nicht möglich
\item
  \textbf{Komplexe Verknüpfungen:} Die Zusammenführung verschiedener
  Datenquellen erfordert effiziente Join-Operationen
\item
  \textbf{Datenschutz:} Strenge Anforderungen an den Umgang mit
  sensiblen Unternehmensdaten
\end{itemize}

\subsection{Zugangsformen im FDZ}\label{zugangsformen-im-fdz}

Das BTP ist über verschiedene Zugangsformen verfügbar:

\begin{itemize}
\tightlist
\item
  \textbf{Gastwissenschaftsarbeitsplatz (GWAP):} On-Site-Zugang mit
  vollständigen Daten, auswertbar mit SAS oder Stata. Bedingt durch die
  Speicherplatzgröße kann die Nutzung einer Stichprobe erforderlich
  sein.
\item
  \textbf{Kontrollierte Datenfernverarbeitung (KDFV):} Remote-Zugang mit
  Ergebnisfreigabe über KDFV
\item
  Die Gemeindekennziffer für Bayern steht am GWAP lediglich
  pseudonymisiert zur Verfügung
\end{itemize}

\subsection{BTP als
Benchmark-Datensatz}\label{btp-als-benchmark-datensatz}

Aufgrund seiner Größe (\textgreater66 Millionen Einheiten), Komplexität
(\textgreater2.700 Variablen) und realen Datenstruktur eignet sich das
BTP hervorragend als Benchmark für die Evaluation verschiedener
Auswertungsmethoden. Die Herausforderungen des BTP sind repräsentativ
für viele Forschungsprojekte im FDZ:

\begin{itemize}
\tightlist
\item
  \textbf{Datengröße:} Der Datensatz übersteigt oft die verfügbare
  RAM-Kapazität
\item
  \textbf{Komplexe Joins:} Verknüpfung verschiedener Steuerstatistiken
  erforderlich\\
\item
  \textbf{Aggregationen:} Über verschiedene Dimensionen (Zeit, Region,
  Wirtschaftszweig)
\item
  \textbf{Filteroperationen:} Selektion relevanter Unternehmensgruppen
\item
  \textbf{Zeitreihenanalysen:} Panel-Struktur über mehrere Jahre
\end{itemize}

Diese realitätsnahen Szenarien ermöglichen es, die Praxistauglichkeit
verschiedener Auswertungsmethoden unter den tatsächlichen Bedingungen
eines Forschungsdatenzentrums zu bewerten.


\printbibliography



\end{document}
